\section{Actual coefficient packets and the coprime resultant test}
\label{cg19-section}

This section returns to the original block. It disproves a proposed
bound based on one short primitive coefficient vector, using actual
different owners, and then proves the correct two-polynomial
certificate with its full signed accounting. The second polynomial
must be coprime to the boundary polynomial; constructing it with a
useful degree and height remains a separate problem.

Let \(n>324\) be prime, \(B=n^4\), and
\(\mathcal O=\mathbb Z[\zeta]\), where \(\zeta^2-\zeta+1=0\) with
positive imaginary part. Put
\[
 P(A+C\zeta)=AC(A+C),\qquad
 w_k=3k+\zeta,\quad z_k=w_k^n,\quad T_k=P(z_k)
 \quad(B\le k<2B).
\]
Here and throughout the block indices are integers. These three
raw arms are positive: if \(\theta=\arg w_k\), then
\(0<n\theta<n/(3k)\le1/3\), so the first coordinate is
\[
 |w_k|^n\left(\cos(n\theta)-\frac{\sin(n\theta)}{\sqrt3}\right)>0,
\]
and the second is positive as well.

The powers are primitive. For a prime \(p\ne3\), the algebra
\(\mathcal O/p\mathcal O\) is reduced because the discriminant is
\(-3\). Thus \(w_k^n=0\) in this algebra would imply \(w_k=0\),
contrary to its second coordinate being one. At \(p=3\), the norm
\(9k^2+3k+1\) is a unit, so again the two coordinates of the power
cannot vanish. Hence the three arms are pairwise coprime. A prime
dividing one arm has nonzero norm at the power and at its root.

A full assigned \(T\)-packet means a finite set of distinct rational
primes \(p\), an actual block owner \(k(p)\) for each, and its full
positive depth \(e_p=v_p(T_{k(p)})\). For a degree limit \(d\) its
evaluation lattice is
\[
 \mathcal L=\{G\in\mathbb Z[X]_{\le d}:
     G(k(p))\equiv0\pmod{p^{e_p}}\text{ for every packet prime}\}.
\]
The statement being tested is that one nonzero primitive vector of
small coefficient height in this lattice bounds its index by a fixed
power of that height, uniformly in \(n\).

\begin{proposition}[Literal primitive boundary polynomial]
\label{cg19-content}
Set
\[
 P_n(a)=P((a+\zeta)^n),\qquad T_n(X)=P_n(3X),\qquad
 F_n(X)=T_n(X)/3.
\]
The content of \(T_n\) is exactly three. Consequently \(F_n\) is
a nonzero primitive integer polynomial of degree \(d=3n-1\), and
\begin{equation}
 \|F_n\|_1\le64^n/3\le64^n.
 \label{cg19-height}
\end{equation}
Moreover there is an \(R_n\in\mathbb Z[a]\) such that
\begin{equation}
 P_n(a)=na+a^2R_n(a).
 \label{cg19-taylor}
\end{equation}
\end{proposition}
\begin{proof}
The sixth roots of unity have coordinate pairs
\[
 (1,0),(0,1),(-1,1),(-1,0),(0,-1),(1,-1).
\]
Since \(n\equiv1\) or \(5\pmod6\), one has \(P_n(0)=0\) and
\(P_n'(0)=n\). In the first case \(\zeta^n=\zeta\); the zero
first arm has derivative \(n\). In the second case
\(\zeta^n=1-\zeta\); the zero sum arm has derivative \(-n\) and
the product of the other two arms is \(-1\). This proves
\eqref{cg19-taylor}. All coefficients of \(T_n\) are divisible
by three, and its linear coefficient is \(3n\).

Frobenius in \((\mathcal O/n\mathcal O)[X]\) gives
\begin{equation}
 T_n(X)\equiv
 \begin{cases}
  (3X)^n((3X)^n+1)&n\equiv1\pmod6,\\
  -(3X)^n((3X)^n+1)&n\equiv5\pmod6
 \end{cases}\pmod n.
 \label{cg19-frobenius}
\end{equation}
This is a nonzero polynomial over \(\mathbb F_n\). The content
divides \(3n\) and is not divisible by \(n\), so it is exactly
three.

Writing \(w_X^n=A_n(X)+C_n(X)\zeta\), the degrees of
\(A_n,C_n,A_n+C_n\) are \(n,n-1,n\), with leading coefficients
\(3^n,n3^{n-1},3^n\). Thus \(\deg F_n=3n-1\).
Every coordinate and the coordinate sum of \(\zeta^j\) has absolute
value at most one. The binomial expansion consequently bounds each
of the three coefficient \(\ell^1\)-norms by
\(\sum_j\binom nj3^{n-j}=4^n\). Submultiplicativity, followed by
division by three, proves \eqref{cg19-height}.
\end{proof}

\begin{proposition}[Two different actual positive-depth owners]
\label{cg19-owners}
For \(k_0=B\),
\[
 v_3(T_{k_0})=1,\qquad v_n(T_{k_0})=5,\qquad v_5(T_{k_0})=0.
\]
There is an integer \(k_1\in[B,2B)\) with
\[
 k_1\equiv5^4\pmod{5^5},\qquad k_1\equiv1\pmod n.
\]
Every such \(k_1\) satisfies
\[
 v_5(T_{k_1})=4,\qquad v_n(T_{k_1})=0.
\]
Thus the two actual distinct owners supply the positive signed
contributions \(2\log n\) and \(\log5\), respectively.
\end{proposition}
\begin{proof}
Since \(3\nmid nB\), equation~\eqref{cg19-taylor} at \(a=3B\)
gives \(v_3(T_B)=1\). At the prime \(n\), use \(a=3n^4\).
At \(a=0\) exactly one arm vanishes, with linear term \(\pm na\)
of valuation five. Every term of degrees two through \(n-1\)
has valuation at least \(1+4\cdot2=9\), since the corresponding
binomial coefficient is divisible by \(n\); the final term has
valuation \(4n\). The other two arms remain units modulo \(n\).
This proves the full depth five.

Modulo five, \(B=n^4=1\), so \(w_B=3+\zeta\) in \(\mathbb F_{25}\).
Its norm is \(13=3\pmod5\), whence \((3+\zeta)^6=3\), and
its fifth power is its conjugate \(4-\zeta\).
As \(n\equiv1\) or \(5\pmod6\), the power \(w_B^n\) is a nonzero
scalar multiple of one of these two elements. Their \(P\)-values
are \(12,-12\), both nonzero modulo five.

CRT supplies the stated congruences with modulus \(3125n\).
The least solution at least \(B\) is less than \(B+3125n<2B\),
because \(n^3>3125\). Its five-adic depth as an index is four.
Equation~\eqref{cg19-taylor} then has linear term \(3nk_1\)
of five-adic depth four and every later term of depth at least
eight. Finally \eqref{cg19-frobenius} at \(k_1\equiv1\pmod n\)
gives \(T_{k_1}\equiv\pm12\pmod n\), a unit.
The owners differ modulo \(n\), so they are distinct.
\end{proof}

These primes are small relative to the moving far-tail cutoffs.
The construction is not a counterexample to a far-tail estimate,
and makes no claim that these particular owners lie in the
multiplicatively independent subdomain.

\begin{theorem}[Exact index and the complete selected ledger]
\label{cg19-index}
Let \(m_0=F_n(B)=T_B/3\). Assign every prime dividing \(m_0\) to
\(k_0\), with its full depth in \(T_B\), and assign the additional
prime five to \(k_1\), with depth four. Then all labels are norm
units at their actual owners, and
\begin{equation}
 M=625m_0,\qquad [\mathbb Z^{d+1}:\mathcal L]=M,\qquad
 F_n\in\mathcal L.
 \label{cg19-index-eq}
\end{equation}
This is not a common-hit rectangle: \(n\) hits only the first of
these owners, and five only the second. For the complete
\(J(T)=\log T-3\log\operatorname{rad}T\), its selected ledger is
\begin{equation}
 \sum_{p\mid m_0}(v_p(m_0)-3)\log p+\log5
 =J(T_B)+2\log3+\log5.
 \label{cg19-ledger}
\end{equation}
\end{theorem}
\begin{proof}
Proposition~\ref{cg19-owners} gives \(\gcd(m_0,15)=1\).
Primitivity in the setup proves the norm-unit assertion.
The additive homomorphism
\[
 \mathbb Z[X]_{\le d}\longrightarrow
 \left(\prod_{p\mid m_0}\mathbb Z/p^{v_p(m_0)}\mathbb Z\right)
 \times\mathbb Z/625\mathbb Z
\]
given by evaluation at the assigned owners is surjective:
constant polynomials already realize all residue tuples by CRT.
Its kernel is \(\mathcal L\), proving the index assertion.
The polynomial \(F_n\) satisfies each first-owner congruence by
its value \(m_0\), and at \(k_1\) it has exact five-adic depth
four because division by three does not affect that depth.
Its coefficient vector is primitive by
Proposition~\ref{cg19-content}, and its value \(F_n(B)=m_0\)
is positive. Thus this is neither the zero polynomial nor a
vanishing integer evaluation used in an additive selector kernel.

The only prime removed from \(T_B\) is three at depth one, whose
contribution to \(J(T_B)\) is \(-2\log3\). The added five-adic
depth four contributes \(\log5\). This proves
\eqref{cg19-ledger}, with every depth-one and depth-two term
retained. Neither its sign nor its asymptotic size is asserted.
\end{proof}

\begin{theorem}[Failure of a fixed-power short-vector bound]
\label{cg19-obstruction}
For \(H_n=\max(2,\|F_n\|_1)\), every fixed \(C,C_0>0\) satisfies
\[
 [\mathbb Z^{d+1}:\mathcal L]>C_0H_n^C
\]
for all sufficiently large prime \(n\) in the construction above.
\end{theorem}
\begin{proof}
Put \(r=|3B+\zeta|\ge3B\) and \(u=n\arg(3B+\zeta)\in(0,1/3)\).
The inequalities \(\cos u\ge1-u^2/2\) and \(\sin u\le u\)
give
\[
 A_B=r^n(\cos u-\sin u/\sqrt3)>r^n/2.
\]
The second coordinate \(C_B\) is a positive integer, and
\(A_B+C_B>A_B\). Therefore
\[
 T_B>(3B)^{2n}/4,\qquad M>(625/12)(3B)^{2n}.
\]
Together with \eqref{cg19-height}, this yields
\[
 \log M>8n\log n+2n\log3+\log(625/12),\qquad
 \log H_n\le n\log64.
\]
The ratio tends to infinity along primes, proving the assertion.
\end{proof}

Only bounds with an absolute fixed exponent based on one primitive
short coefficient vector are refuted. The dimension grows with \(n\).
A dimension-dependent bound, a bound paying the evaluation height
\(B\), and Hadamard applied to a full independent basis are not
refuted. In particular, one short primitive lattice vector is
insufficient evidence for a small covolume.

\begin{theorem}[A coprime second polynomial certificate]
\label{cg19-resultant}
Let \(F,G\in\mathbb Z[X]\setminus\{0\}\), of actual degrees
\(d,e\ge1\), be coprime over \(\mathbb Q\).
For a finite set of distinct primes, positive integers \(h_p\),
and arbitrary integer owners \(a_p\), suppose
\[
 p^{h_p}\mid F(a_p),\qquad p^{h_p}\mid G(a_p).
\]
Then
\begin{equation}
 \prod_p p^{h_p}\mid\operatorname{Res}(F,G),\qquad
 \sum_p h_p\log p\le e\log\|F\|_2+d\log\|G\|_2.
 \label{cg19-resultant-bound}
\end{equation}
No monicity or invertibility of a leading coefficient is required.
The empty packet is allowed, with modulus product one.
\end{theorem}
\begin{proof}
The integer Sylvester determinant is nonzero by coprimality.
View its matrix as the coefficient map
\((U,V)\mapsto UF+VG\), for \(\deg U<e\) and \(\deg V<d\).
Applying its integer adjugate to the constant basis vector gives
\(U,V\in\mathbb Z[X]\) satisfying
\[
 UF+VG=\pm\operatorname{Res}(F,G).
\]
Evaluation at each assigned owner proves divisibility by the full
prime power \(p^{h_p}\). Multiplying these pairwise coprime
divisibilities proves the first assertion. This argument never
passes to a field root or divides by a leading coefficient.

The Sylvester determinant has \(e\) shifted coefficient rows of
\(F\) and \(d\) of \(G\). Their Euclidean lengths are respectively
\(\|F\|_2,\|G\|_2\). Hadamard's inequality gives
\[
 |\operatorname{Res}(F,G)|\le\|F\|_2^e\|G\|_2^d.
\]
Taking logarithms of the nonzero integer multiple proves the
second assertion, including all prime-power depths.
\end{proof}

For a nonzero constant \(G=c\), the direct identity
\(\operatorname{Res}(F,c)=c^d\) covers the degree-zero case.
Each packet modulus must then divide \(c\); it is no useful short
certificate unless the modulus is already small.
Theorem~\ref{cg19-resultant} genuinely allows different primes
to have different owners. It requires one coprime second
polynomial, rather than a complete lattice basis, but its saving
depends on constructing that polynomial with suitable height.

\begin{proposition}[Automatic owner polynomial and the exact correction at three]
\label{cg19-automatic}
First take a nonempty full assigned \(T\)-packet whose primes
are different from three, and let \(\mathcal K\) be its set of
distinct owners, \(s=|\mathcal K|\). Then
\[
 G_{\mathcal K}(X)=\prod_{k\in\mathcal K}(X-k)
\]
is a coprime second certificate for \(F_n\), with
\begin{align}
 |\operatorname{Res}(F_n,G_{\mathcal K})|
   &=\prod_{k\in\mathcal K}F_n(k), \label{cg19-auto-res}\\
 \|G_{\mathcal K}\|_2\le\|G_{\mathcal K}\|_1
   &=\prod_{k\in\mathcal K}(1+k)\le(2B)^s. \label{cg19-auto-norm}
\end{align}
Thus its Sylvester estimate gives only
\[
 \log M\le s\log\|F_n\|_2+(3n-1)s\log(2B).
\]

For an arbitrary nonempty full \(T\)-packet including three, let
\(k_3\) be its assigned owner. Its exact depth is
\begin{equation}
 e_3=1+v_3(k_3).
 \label{cg19-three-depth}
\end{equation}
Form the adjusted \(F_n\)-packet with \(h_p=e_p\) for \(p\ne3\)
and \(h_3=e_3-1\), omitting three when \(h_3=0\). Its product
\(M_F\) and the original product \(M_T\) satisfy
\begin{equation}
 M_T=3M_F.
 \label{cg19-three-product}
\end{equation}
The same owner polynomial is a certificate for the adjusted
packet, even when it is empty, and the logarithmic bound gains
exactly one \(\log3\) for \(M_T\).
For the respective positive-depth signed supports,
\begin{equation}
 J_T=
 \begin{cases}
  J_F-2\log3,&e_3=1,\\
  J_F+\log3,&e_3\ge2.
 \end{cases}
 \label{cg19-three-signed}
\end{equation}
\end{proposition}
\begin{proof}
At every owner \(F_n(k)=T_k/3>0\); hence none is a root of \(F_n\),
so the owner polynomial is coprime to it. It vanishes as an
integer at every owner, proving all adjusted congruences.
The product formula for the resultant with its monic linear
factors proves \eqref{cg19-auto-res}. Since every \(k>0\), the
absolute elementary symmetric coefficients sum to
\(\prod(1+k)\); this proves \eqref{cg19-auto-norm}.
Theorem~\ref{cg19-resultant} gives the displayed bound.

Division by three is the reason the first assertion excluded
that prime. For a general block owner, substitute \(a=3k_3\)
into \eqref{cg19-taylor}. The linear term has valuation
\(1+v_3(k_3)\), because \(3\nmid n\); every later term has
valuation at least \(2+2v_3(k_3)\). This proves
\eqref{cg19-three-depth}, with no assumption that \(k_3\) is a
three-adic unit. Division of the value by three reduces this
depth by exactly one, proving \eqref{cg19-three-product}.

If \(e_3=1\), then
\(\operatorname{rad}M_T=3\operatorname{rad}M_F\); if \(e_3\ge2\)
their radicals are equal. Subtracting three times their radical
logs from \(\log M_T=\log M_F+\log3\) proves
\eqref{cg19-three-signed}. Thus even the depth-one and depth-two
negative contributions are translated exactly.
\end{proof}

\begin{corollary}[Signed degree-height test with the original radical]
\label{cg19-signed-certificate}
Let \(G\in\mathbb Z[X]\setminus\{0\}\) be coprime to \(F_n\)
over \(\mathbb Q\) and satisfy the
adjusted \(F_n\)-packet congruences. For the adjusted packet,
\[
 J_F\le
 \deg(G)\log\|F_n\|_2+(3n-1)\log\|G\|_2
       -3\log\operatorname{rad}M_F.
\]
Equivalently, for its original full \(T\)-packet,
\begin{equation}
 J_T\le
 \deg(G)\log\|F_n\|_2+(3n-1)\log\|G\|_2
       +\delta_3\log3-3\log\operatorname{rad}M_T,
 \label{cg19-full-signed}
\end{equation}
where \(\delta_3\) is one if three is present and zero otherwise.
\end{corollary}
\begin{proof}
For positive degree apply Theorem~\ref{cg19-resultant} to the
adjusted product; for constant \(G=c\ne0\), use the preceding
direct case \(\operatorname{Res}(F_n,c)=c^{3n-1}\).
Then use \(J=\log M-3\log\operatorname{rad}M\).
If three is present, use \eqref{cg19-three-product}; otherwise
the products coincide. This gives \eqref{cg19-full-signed}
without discarding a negative term.
\end{proof}

\paragraph{What the automatic certificate does and does not prove.}
Equation~\eqref{cg19-auto-res} recovers the bound obtained by
assigning each prime power to its owner and multiplying the
owner values. In the actual two-owner example it is
\((X-B)(X-k_1)\), with resultant \(F_n(B)F_n(k_1)\). This pays
the old owner height and is consistent with
Theorem~\ref{cg19-obstruction}. It does not itself supply the
negative prime-depth compensation needed for a signed saving.

A genuine new gain requires a coprime certificate whose
\[
 \deg(G)\log\|F_n\|_2+(3n-1)\log\|G\|_2
\]
improves the relevant full signed budget. A small lattice
vector may simply be a multiple of \(F_n\), while the automatic
product of owner factors gives \eqref{cg19-auto-res}. None of
the results cited here constructs the required shorter coprime
certificate on the original independent domain, and no
dependence on the packet modulus is hidden in its claimed height.

\paragraph{Ordinary and finite scope.}
The arguments in this section are complete ordinary proofs.
The separate exact coefficient-packet replay checks three prime
indices \(331,337,347\), six actual powers, coefficient content
and degree, independent polynomial evaluation, both exact marked
depths and CRT indices, and the automatic resultant divisibility.
It neither fully factors the first-owner boundary nor computes
its full signed ledger, and finite cases are not the proof of
the asymptotic obstruction. No Lean formalization of this
resultant and lattice chain is claimed here.

The positive fixed-exponent selector families of
Sections~\ref{np19-section} and~\ref{cu19-section} are separate.
These results do not exclude useful small intersections or
other nonlinear maps, do not control the original far signed
tail, and do not close its exceptional set, the dependent-root
complement, or the reduction from this special block to every
primitive ABC triple.
